%% =====================================================================
%%  B1-T-16-04 -- what B1 contributes to "The Unargue"
%%  -------------------------------------------------------------------
%%  LAYER A: APPEND-ONLY. Written once at the entry's write, then left
%%  alone. If the source has more to say later, add a bullet -- do not
%%  rewrite. Material found ONLY here goes in the `unique` slot with
%%  @unique set to yes, which makes it undeletable (rule R10).
%% =====================================================================
%% @kind: source
%% @id: B1-T-16-04
%% @book: B1
%% @topic: T-16-04
%% @locator: Tactic 11 (ch. XII), pp. 125--152
%% @status: distilled
%% @unique: yes
%% @integrated: yes
%% @updated: 2026-09-19

\dossier{B1}{T-16-04}{\bookshort{B1}}{Tactic 11 (ch. XII), pp. 125--152}

\begin{distilled}
  \item An argument yields an angry adversary more convinced than ever; Unargue is the opposite --- an ex-adversary ready to let you have your way.
  \item Step 1, the soulmate tactic: dwell on real areas of agreement, no pretence --- anyone over an IQ of fifty spots a phony and resents it.
  \item Step 2: listen, listen, listen --- hear out the objection, sort token misgivings from the one or two major ones; the target feels cared-for instead of ignored.
  \item Step 3: agree with the feelings, not the case --- I know how you feel, plus a self-deprecating laugh, massages the ego aside without conceding substance.
  \item Step 4: point out areas of agreement and appeal to common values --- we both want this campaign to make you money.
  \item Step 5: state your case through the nine-out-of-ten device --- you'd usually be right, but this situation is unusual; handle the major objections and momentum routs the rest, like the loss of a major battle.
  \item Step 6: save face throughout --- concede a minor point to start a pattern of concession, offer a compromise so the target can yield with head held high.
  \item The chapter's admission is in its title: manipulate a person against his will and make him like it.
  \item Token-objection doctrine: the stated objection is often a surrogate probing your handling; the real one is unspoken.
\end{distilled}
\begin{unique}
  \item The six-step Unargue sequence itself, the soulmate tactic, the nine-out-of-ten device, and the token-objection doctrine are unique to B1 in this corpus.
\end{unique}
\begin{terms}
  \item \emph{the Unargue} --- Sparkman's name for his six-step disagreement technique that reverses the argument's outcome.
  \item \emph{yes-no state} --- B1's term for the target's paralysis between wanting and refusing, the trigger for the pressure family.
  \item \emph{token objection} --- a surrogate objection offered to test how the operator handles resistance; the real objection is held back.
\end{terms}
